% !TeX root = ../main.tex

% ==================================================
% D.5. Index alignment와 이항 연산
% ==================================================

\subsection{Index alignment와 이항 연산}

\subsubsection{이전 학습과의 연결}

\zsh{D.4}에서는 Pandas의 데이터를 선택하는 두 가지 좌표계를 구분했다.

\begin{dhksmry}
    - 레이블(label): Index가 관리하는 식별값

    - 위치(position): 현재 데이터에서 원소가 놓인 순서
\end{dhksmry}

\zsh{df.loc[...]}은 레이블을 기준으로, \zsh{df.iloc[...]}은 위치를 기준으로 데이터를 선택했다.

그런데 Pandas에서 Index의 역할은 데이터 선택에 그치지 않는다.

두 Series를 더하거나, DataFrame에 Series를 할당하거나, 

서로 다른 DataFrame을 연산할 때도 Index가 결과에 개입한다.

다음 코드를 살펴보자.

\begin{dhkcode}
import pandas as pd

a = pd.Series(
    [10, 20],
    index=["A", "B"],
)

b = pd.Series(
    [1, 2],
    index=["B", "A"],
)

print(a + b)
\end{dhkcode}

예상 출력:

\begin{dhkzsh}
A    12
B    21
dtype: int64
\end{dhkzsh}

두 Series의 \zsh{shape}은 모두 \zsh{(2,)}다.

하지만 위치별로 더한 \zsh{[11, 22]}가 아니라 \zsh{[12, 21]}이 반환되었다.

왜 이런 결과가 발생할까?

Pandas는 두 Series를 더할 때 단순히 같은 위치의 원소를 대응시키는 것이 아니라, Index 레이블이 같은 원소끼리 대응시킨 다음 연산하기 때문이다.

이 과정을 \zsh{Index alignment(인덱스 정렬 또는 레이블 정합)}라고 한다.

여기서 alignment를 단순한 행 정렬(sort)과 혼동하면 안 된다.

Alignment의 핵심은 데이터를 어떤 순서로 나열하는가가 아니라, 어떤 원소와 어떤 원소가 서로 대응하는가를 결정하는 데 있다.

\subsubsection{이번 학습의 핵심 질문}

이번 \zsh{D.5}에서는 다음 내용을 분석할 것이다.

\begin{dhksmry}
Index alignment는 정확히 어떤 연산이며, 언제 자동으로 발생하는가?

두 Series의 레이블 순서가 다르거나 일부 레이블이 없으면 어떻게 되는가?

DataFrame끼리 연산하면 행 Index와 열 Index가 각각 어떤 역할을 하는가?

DataFrame과 Series를 연산할 때는 어느 축을 기준으로 대응시키는가?

산술 연산, 비교 연산, 할당은 모두 같은 정렬 규칙을 따르는가?

중복 Index가 있으면 결과의 행 수와 값이 어떻게 달라지는가?

NumPy 배열로 변환하면 왜 레이블 정합성이 사라지는가?
\end{dhksmry}

이번 학습의 가장 중요한 목표는 다음 판단을 할 수 있게 되는 것이다.

\begin{dhksmry}[centered]
연산을 수행하기 전에 두 객체의 값뿐 아니라 Index를 확인하여, 어떤 

원소끼리 계산되고 어떤 결과가 결측으로 남을지 예측할 수 있어야 한다.
\end{dhksmry}

\subsubsection{Index alignment의 정확한 정의}

Index alignment는 두 Pandas 객체의 연산에서 레이블을 기준으로 대응 관계를 구성하는 과정이다.

Series 두 개가 있다고 가정해 보자.

\begin{dhkcode}
a = pd.Series([10, 20], index=["A", "B"])
b = pd.Series([1, 2], index=["B", "A"])
\end{dhkcode}

각 Series는 다음과 같은 대응 관계를 갖고 있다.

\begin{dhkmatrix}[delim=none,style=plain]
cols:: ,Series a, , ,Series b
data :: [[A,10, <->, B,1],[B,20,레이블 대응,A,2]]
\end{dhkmatrix}

\begin{dhkmatrix}[delim=none,style=plain]
title::Alignment 이후의 계산 관계
data::[[레이블 A,,10+2, ,12],[레이블 B,,20+1, ,21]]
\end{dhkmatrix}

두 Series의 Index 레이블이 유일하다고 가정하면, 기본 산술 연산의 대응 관계를 다음과 같이 표현할 수 있다.

\[
(a+b)[k]=a[k]+b[k]
\]

여기서 \(k\)는 Index 레이블이다.

두 Series의 레이블 집합이 다르다면 기본 산술 연산은 일반적으로 양쪽 레이블의 합집합(union)을 기준으로 결과를 구성한다.

\[
I_{\text{result}}=I_a\cup I_b
\]

다만 이것은 유일한 Index를 갖는 Series의 기본 산술 연산을 설명하는 모델이다.

중복 레이블, 비교 연산, 명시적인 정렬 옵션 등에서는 세부 규칙이 달라질 수 있다.

또한 이 수식은 계산할 값의 대응 관계를 나타내는 것이지, Pandas 내부에서 반드시 레이블별 Python 반복문을 실행한다는 뜻은 아니다.

Pandas는 Index에 맞춰 데이터의 위치 관계를 구성한 다음, 배열 연산을 활용해 계산할 수 있다.

\subsubsection{Series간 산술 연산 - 레이블 순서가 다른 경우}

먼저 \zsh{D.5.1}의 예제를 객체 관점에서 다시 분석해 보자.

\begin{dhknote}
\begin{dhksmry}[centered]
독립 실행 코드 1
\end{dhksmry}
\end{dhknote}

\begin{dhkcode}
import numpy as np
import pandas as pd

a = pd.Series(
    [10, 20],
    index=["A", "B"],
)

b = pd.Series(
    [1, 2],
    index=["B", "A"],
)

result = a + b

print(result)
print(result.shape)
print(result.dtype)
print(a.index.equals(b.index))

print(a.to_numpy() + b.to_numpy())
\end{dhkcode}

예상 출력:

\begin{dhkzsh}
A    12
B    21
dtype: int64
(2,)
int64
False
[11 22]
\end{dhkzsh}

\begin{dhknote}
\begin{dhksmry}[centered]
실행 과정 분석
\end{dhksmry}
\end{dhknote}

\begin{dhknote}
\begin{dhksmry}[centered]
1단계 - Series 객체 생성
\end{dhksmry}
\end{dhknote}

\zsh{a}와 \zsh{b}는 서로 다른 Series 인스턴스를 참조한다.

두 객체 모두 \zsh{shape == (2,)}이고 \zsh{dtype == int64} 이지만, Index 레이블의 순서가 다르다.

\begin{dhknote}
\begin{dhksmry}[centered]
2단계 - a + b 평가
\end{dhksmry}
\end{dhknote}

Python은 \zsh{+} 연산에 대응하는 객체의 산술 연산 프로토콜을 사용한다.

Pandas Series는 이러한 연산을 구현하면서 Index alignment를 수행한다.

따라서 \zsh{a}의 \zsh{"A"}에 대응하는 값 10과 \zsh{b}의 \zsh{"A"}에 대응하는 값 2가 더해진다.

\begin{dhknote}
\begin{dhksmry}[centered]
3단계 - 결과 객체 반환
\end{dhksmry}
\end{dhknote}

연산 결과를 담은 새로운 Series가 반환되고, \zsh{result}라는 이름에 바인딩된다.

원본 \zsh{a}, \zsh{b}의 값과 Index는 변경되지 않는다.

\begin{dhknote}
\begin{dhksmry}[centered]
4단계 - NumPy 배열과 비교
\end{dhksmry}
\end{dhknote}

\begin{dhkcode}
a.to_numpy() + b.to_numpy()
\end{dhkcode}

이번에는 Series가 아니라 NumPy 배열끼리 더한다.

NumPy 배열에는 Pandas의 Index 레이블이 없으므로 같은 위치의 값이 대응된다.

따라서 결과가 \zsh{[11, 22]}로 달라지는 것이다.

이것은 \zsh{D.1}에서 배운 NumPy와 Pandas의 가장 중요한 차이를 실제 연산으로 확인한 사례다.

Pandas의 Series끼리 연산하는 것과 Series를 NumPy 배열로 변환한 뒤 연산하는 것은 일반적으로 동등하지 않다.

\subsubsection{두 Index가 완전히 일치하지 않으면 어떻게 될까?}

지금까지는 두 Series의 레이블 집합이 같았다.

이번에는 일부 레이블만 일치하는 경우를 살펴보자.

\begin{dhknote}
\begin{dhksmry}[centered]
독립 실행 코드 2 - 레이블 합집합과 결측값
\end{dhksmry}
\end{dhknote}

\begin{dhkcode}
import pandas as pd

a = pd.Series(
    [10, 20],
    index=["A", "B"],
    dtype="int64",
)

b = pd.Series(
    [1, 2],
    index=["B", "C"],
    dtype="int64",
)

result = a + b

print(result)
print(result.shape)
print(result.dtype)
\end{dhkcode}

예상 출력:

\begin{dhkzsh}
A     NaN
B    21.0
C     NaN
dtype: float64
(3,)
float64
\end{dhkzsh}

왜 결과가 세 개의 원소를 갖게 되었을까?

두 Series의 레이블을 비교해 보자.

\begin{dhktable}
|레이블|	a의 값|	b의 값|	연산 결과|
|:---|:---|:---|:---|
|A|	10|	없음	|NaN|
|B|	20|	1	|21.0|
|C|	없음|	2|	NaN|
\end{dhktable}

기본 산술 연산은 두 Index의 합집합인 \zsh{["A", "B", "C"]}를 결과 레이블로 구성한다.

그런데 \zsh{"A"}에는 \zsh{b}의 값이 없고, \zsh{"C"}에는 \zsh{a}의 값이 없다.

이 경우 계산에 필요한 피연산자가 존재하지 않으므로 결측값이 발생한다.

\begin{dhknote}
\begin{dhksmry}[centered]
왜 int64가 float64로 바뀌었을까?
\end{dhksmry}
\end{dhknote}

\zsh{D.3}에서 배운 내용을 연결해 보자.

일반적인 NumPy 기반 \zsh{int64}는 \zsh{NaN}을 정수 결측값으로 표현할 수 없다.

따라서 결측이 발생한 결과는 \zsh{float64}로 변환될 수 있다.

하지만 Nullable 정수형이라면 결과가 달라진다.

\begin{dhkcode}
import pandas as pd

a = pd.Series(
    [10, 20],
    index=["A", "B"],
    dtype="Int64",
)

b = pd.Series(
    [1, 2],
    index=["B", "C"],
    dtype="Int64",
)

result = a + b

print(result)
\end{dhkcode}

예상 출력:

\begin{dhkzsh}
A    <NA>
B      21
C    <NA>
dtype: Int64
\end{dhkzsh}

이번에는 \zsh{Int64} 자료형이 결측을 표현할 수 있기 때문에 결과도 Nullable 정수형으로 유지된다.

따라서 다음 두 가지를 구분해야 한다.

\begin{dhksmry}
    - Alignment는 어떤 레이블끼리 값을 대응시킬지 결정한다.

    - dtype과 연산 규칙은 대응되는 값이 없을 때 결측을 어떻게 표현할지
      결정한다.
\end{dhksmry}

Index alignment 자체가 무조건 \zsh{float64}를 만든다는 뜻은 아니다.

\subsubsection{\zsh{fill_value}는 결측값을 어떻게 처리할까?}

기본 산술 연산에서는 대응되는 값이 없으면 결측이 발생할 수 있다.

하지만 Pandas는 산술 메서드에 \zsh{fill_value} 인수를 제공한다.

\begin{dhknote}
\begin{dhksmry}[centered]
독립 실행 코드 3 - add()와 fill_value
\end{dhksmry}
\end{dhknote}

\begin{dhkcode}
import pandas as pd

a = pd.Series(
    [10, 20],
    index=["A", "B"],
)

b = pd.Series(
    [1, 2],
    index=["B", "C"],
)

result = a.add(
    b,
    fill_value=0,
)

print(result)
\end{dhkcode}

예상 출력:

\begin{dhkzsh}
A    10.0
B    21.0
C     2.0
dtype: float64
\end{dhkzsh}

\zsh{Series.add()}는 Series의 산술 메서드다.

\begin{dhkcode}
a.add(b, fill_value=0)
\end{dhkcode}

이 표현식에서는 \zsh{a}에 바인딩된 메서드를 호출하고, \zsh{b}와 \zsh{fill_value=0}을 인수로 전달한다.

기본적인 레이블 대응은 그대로 유지하면서, 한쪽 값이 결측인 경우 0을 대신 사용해 계산할 수 있다.

\begin{dhktable}
|레이블|	a|	b|	계산|
|:---|:---|:---|:---|
|A|	10|	없음|	10 + 0|
|B|	20|	1	|20 + 1|
|C|	없음|	2|	0 + 2|
\end{dhktable}

그런데 여기에는 중요한 예외가 있다.

\begin{dhknote}
\begin{dhksmry}[centered]
양쪽 값이 모두 결측이라면?
\end{dhksmry}
\end{dhknote}

\begin{dhkcode}
import numpy as np
import pandas as pd

a = pd.Series(
    [10, np.nan],
    index=["A", "B"],
)

b = pd.Series(
    [1, np.nan],
    index=["A", "C"],
)

print(a.add(b, fill_value=0))
\end{dhkcode}

예상 출력:

\begin{dhkzsh}
A    11.0
B     NaN
C     NaN
dtype: float64
\end{dhkzsh}

\zsh{fill_value=0}을 지정했는데도 \zsh{"B"}와 \zsh{"C"}는 \zsh{NaN}이다.

이유는 \zsh{add(fill_value=0)}이 양쪽 모두 결측인 위치까지 무조건 0으로 바꾸는 연산은 아니기 때문이다.

한쪽만 결측인 경우에는 대체값을 사용할 수 있지만, 양쪽 모두 결측인 경우에는 결측 결과가 유지된다.

반면 다음 코드는 다르다.

\begin{dhkcode}
result = (a + b).fillna(0)

print(result)
\end{dhkcode}

예상 출력:

\begin{dhkzsh}
A    11.0
B     0.0
C     0.0
dtype: float64
\end{dhkzsh}

이 경우에는 산술 연산을 먼저 수행한 다음, 결과에 존재하는 결측값을 모두 0으로 바꾼다.

따라서 다음 두 표현식은 일반적으로 동등하지 않다.

\begin{dhkcode}
a.add(b, fill_value=0)

(a + b).fillna(0)
\end{dhkcode}

첫 번째는 연산 과정에서 결측 피연산자를 대체하는 방식이고, 두 번째는 연산 결과의 결측을 대체하는 방식이다.

이 차이는 실험 지표를 합산할 때 특히 중요하다.

관측값이 없어서 계산할 수 없는 결과와 실제 계산 결과가 0인 상황은 서로 다른 의미를 갖기 때문이다.

\subsubsection{명시적 Alignment - \zsh{align()}과 \zsh{reindex()}}

지금까지는 \zsh{+} 연산을 수행하면서 Pandas가 자동으로 Index alignment를 처리했다.

그런데 실제 분석에서는 연산하기 전에 두 객체의 레이블을 명시적으로 맞춰야 할 때가 있다.

이때 사용할 수 있는 메서드가 \zsh{align()}과 \zsh{reindex()}다.

\begin{dhknote}
\begin{dhksmry}[centered]
Series.align()의 구조
\end{dhksmry}
\end{dhknote}

\zsh{align()}은 두 객체의 축을 지정된 방식으로 맞춘 결과를 반환한다.

\begin{dhkcode}
aligned_a, aligned_b = a.align(
    b,
    join="inner",
)
\end{dhkcode}

여기서 \zsh{join}은 어떤 레이블 집합을 사용할지 결정하는 인수다.

\begin{dhktable}
|join	|레이블 선택 기준|
|:---|:---|
|"outer"|	양쪽 레이블의 합집합|
|"inner"|	양쪽 레이블의 교집합|
|"left"|	왼쪽 객체의 레이블|
|"right"|	오른쪽 객체의 레이블|
\end{dhktable}

\begin{dhknote}
\begin{dhksmry}[centered]
독립 실행 코드 4
\end{dhksmry}
\end{dhknote}

\begin{dhkcode}
import pandas as pd

a = pd.Series(
    [10, 20],
    index=["A", "B"],
)

b = pd.Series(
    [1, 2],
    index=["B", "C"],
)

left, right = a.align(
    b,
    join="inner",
)

print(left)
print(right)
print(left + right)
\end{dhkcode}

예상 출력:

\begin{dhkzsh}
B    20
dtype: int64

B    1
dtype: int64

B    21
dtype: int64
\end{dhkzsh}

\zsh{join="inner"}를 지정했기 때문에 두 Series에 공통으로 존재하는 \zsh{"B"}만 유지되었다.

여기서 반환되는 객체 관계를 분석해 보자.

\begin{dhkcode}
left, right = a.align(b, join="inner")
\end{dhkcode}

\zsh{a.align(...)}은 정렬된 두 객체를 담은 튜플을 반환한다.

Python의 언패킹(unpacking)에 의해 첫 번째 결과가 \zsh{left}, 두 번째 결과가 \zsh{right}에 바인딩된다.

원본 \zsh{a}, \zsh{b}는 변경되지 않는다.

\begin{dhknote}
\begin{dhksmry}[centered]
reindex()와의 차이
\end{dhksmry}
\end{dhknote}

\zsh{reindex()}는 하나의 객체를 지정된 레이블 구조에 맞추는 메서드다.

\begin{dhkcode}
import pandas as pd

a = pd.Series(
    [10, 20],
    index=["A", "B"],
)

b = pd.Series(
    [1, 2],
    index=["B", "C"],
)

aligned_b = b.reindex(a.index)

print(aligned_b)
\end{dhkcode}

예상 출력:

\begin{dhkzsh}
A    NaN
B    1.0
dtype: float64
\end{dhkzsh}

\zsh{b.reindex(a.index)}는 \zsh{a.index}를 목표 Index로 사용한다.

따라서 \zsh{"A"}는 결측이 되고, \zsh{"B"}는 기존 값을 유지한다.

\zsh{"C"}는 목표 Index에 없으므로 결과에서 제외된다.

두 메서드의 차이는 다음과 같다.

\begin{dhktable}
|메서드	|주요 목적|	반환|
|:---|:---|:---|
|a.align(b)|	두 객체의 축을 함께 맞춤|	두 객체를 담은 튜플|
|a.reindex(index)|	한 객체를 목표 Index에 맞춤	|재인덱싱된 객체|
\end{dhktable}

여기서 주의할 점은 \zsh{reindex()}가 단순한 순서 변경 메서드가 아니라는 사실이다.

목표 Index에 없는 레이블을 제거하거나 새로운 레이블에 결측값을 생성할 수도 있다.

따라서 \zsh{reindex()}를 사용하기 전에 레이블 집합이 일치하는지 확인하는 것이 중요하다.

\subsubsection{DataFrame 간 연산에서는 두 축이 모두 정렬된다}

Series에는 하나의 Index 축이 있었다.

반면 DataFrame에는 행 Index와 열 Index가 존재한다.

따라서 두 DataFrame을 연산할 때는 행과 열의 레이블이 모두 대응되어야 한다.

\begin{dhknote}
\begin{dhksmry}[centered]
독립 실행 코드 5 - 두 축의 Alignment
\end{dhksmry}
\end{dhknote}

\begin{dhkcode}
import pandas as pd

a = pd.DataFrame(
    {
        "x": [1, 2],
        "y": [10, 20],
    },
    index=["A", "B"],
)

b = pd.DataFrame(
    {
        "y": [100, 200],
        "z": [1000, 2000],
    },
    index=["B", "C"],
)

result = a + b

print(result)
print(result.shape)
\end{dhkcode}

예상 출력:

\begin{dhkzsh}
    x      y   z
A NaN    NaN NaN
B NaN  120.0 NaN
C NaN    NaN NaN

(3, 3)
\end{dhkzsh}

두 DataFrame의 구조를 비교해 보자.

\begin{dhkmatrix}[delim=none,style=cards]
title::DataFrame a
rows:: A,B
cols::x,y
data::[[1,10],[2,20]]
\end{dhkmatrix}

\begin{dhkmatrix}[delim=none,style=cards]
title::DataFrame b
rows:: B,C
cols::y,z
data::[[100,1000],[200,2000]]
\end{dhkmatrix}

\begin{dhksmry}[centered]
공통으로 대응되는 셀
행 B x 열 y
20 + 100 = 120
\end{dhksmry}

두 객체에서 공통으로 존재하는 행 레이블은 \zsh{"B"}다.

공통으로 존재하는 열 레이블은 \zsh{"y"}이다.

따라서 양쪽에 값이 모두 존재하는 셀은 \zsh{("B", "y")} 하나뿐이다.

나머지 셀은 대응되는 피연산자가 없으므로 \zsh{NaN}이 된다.

여기서 중요한 것은 DataFrame끼리의 산술 연산은 행과 열의 위치가 아니라 두 축의 레이블을 기준으로 수행된다는 점이다.

또한 \zsh{a.shape == b.shape == (2,2)}였지만 결과는 \zsh{(3, 3)}이 되었다.

따라서 동일한 \zsh{shape}을 가진 두 DataFrame을 연산한다고 해서 결과의 \zsh{shape}이 반드시 유지되는 것은 아니다.

\subsubsection{DataFrame과 Series의 연산은 어느 축을 기준으로 정렬될까?}

이 부분은 Pandas를 처음 학습할 때 특히 혼동되기 쉽다.

두 DataFrame의 연산에서는 두 축이 모두 정렬되었다.

그렇다면 DataFrame과 Series를 연산하면 Series의 Index가 DataFrame의 행에 대응할까, 열에 대응할까?

기본적으로 DataFrame과 Series의 산술 연산은 Series의 Index를 DataFrame의 열 레이블에 대응시킨다.

\begin{dhknote}
\begin{dhksmry}[centered]
독립 실행 코드 6 - 기본 축과 명시적 축
\end{dhksmry}
\end{dhknote}

\begin{dhkcode}
import pandas as pd

df = pd.DataFrame(
    {
        "label": [1, 0],
        "score": [10, 20],
    },
    index=["doc-A", "doc-B"],
)

bias = pd.Series(
    [100, 200],
    index=["doc-A", "doc-B"],
)

print(df + bias)
\end{dhkcode}

예상 출력:

\begin{dhkzsh}
       doc-A  doc-B  label  score
doc-A    NaN    NaN    NaN    NaN
doc-B    NaN    NaN    NaN    NaN
\end{dhkzsh}

왜 모든 값이 \zsh{NaN}이 되었을까?

\zsh{bias}의 Index는 다음과 같다.

\begin{dhkzsh}
doc-A, doc-B
\end{dhkzsh}

반면 \zsh{df}의 열 레이블은 다음과 같다.

\begin{dhkzsh}
label, score
\end{dhkzsh}

기본 산술 연산은 Series의 Index를 DataFrame의 열 레이블에 대응시키려 한다.

그런데 일치하는 레이블이 없으므로 계산 가능한 값이 존재하지 않는다.

\begin{dhknote}
\begin{dhksmry}[centered]
행 기준으로 정렬하려면?
\end{dhksmry}
\end{dhknote}

\begin{dhkcode}
result = df.add(
    bias,
    axis="index",
)

print(result)
\end{dhkcode}

예상 출력:

\begin{dhkzsh}
       label  score
doc-A    101    110
doc-B    200    220
\end{dhkzsh}

이번에는 \zsh{axis="index"}를 명시했다.

따라서 Series의 Index가 DataFrame의 행 Index에 대응한다.

계산 과정은 다음과 같다.

\begin{dhktable}
|행 레이블|	bias|	label 계산|	score 계산|
|:---|:---|:---|:---|
|doc-A|	100	|1 + 100	|10 + 100|
|doc-B|	200	|0 + 200|	20 + 200|
\end{dhktable}

반대로 열별로 다른 값을 더하려면 Series의 Index를 열 레이블로 구성하면 된다.

\begin{dhkcode}
column_bias = pd.Series({
    "label": 1,
    "score": 2,
})

print(df + column_bias)
\end{dhkcode}

예상 출력:

\begin{dhkzsh}
       label  score
doc-A      2     12
doc-B      1     22
\end{dhkzsh}

이 차이를 정리하면 다음과 같다.

\begin{dhktable}
|연산|	Series의 Index가 대응되는 축|
|:---|:---|
|df + s|	기본적으로 DataFrame의 열|
|df.add(s, axis="columns")	|DataFrame의 열|
|df.add(s, axis="index")	|DataFrame의 행|
\end{dhktable}

NumPy의 브로드캐스팅(broadcasting)과 비교해 보면 차이가 분명해진다.

NumPy에서는 주로 배열의 차원과 길이를 기준으로 브로드캐스팅 가능성을 판단한다.

반면 Pandas에서는 DataFrame과 Series 사이의 산술 연산에 레이블 대응 규칙이 함께 적용된다.

따라서 Series의 길이가 DataFrame의 행 수와 같다는 사실만으로 행 기준 연산이 수행되는 것은 아니다.

\subsubsection{산술 연산과 비교 연산은 같은 규칙을 따를까?}

지금까지 \zsh{+} 연산을 중심으로 살펴보았다.

그런데 \zsh{==} 같은 비교 연산도 항상 같은 방식으로 자동 정렬될까?

그렇지는 않다.

\begin{dhknote}
\begin{dhksmry}[centered]
독립 실행 코드 7 - ==와 .eq() 비교
\end{dhksmry}
\end{dhknote}

\begin{dhkcode}
import pandas as pd

a = pd.Series(
    [10, 20],
    index=["A", "B"],
)

b = pd.Series(
    [20, 10],
    index=["B", "A"],
)

try:
    print(a == b)
except ValueError as e:
    print(type(e).__name__)

print(a.eq(b))
\end{dhkcode}

예상 출력:

\begin{dhkzsh}
ValueError
A    True
B    True
dtype: bool
\end{dhkzsh}

두 Series는 같은 레이블 집합을 갖고 있지만 순서가 다르다.

\zsh{a == b}는 Pandas의 Series 비교 연산 규칙에 따라, 두 객체의 레이블이 동일한 순서로 배치되어 있지 않으면 \zsh{ValueError}를 발생시킨다.

반면 \zsh{a.eq(b)}는 레이블을 정렬하여 비교할 수 있다.

여기서 중요한 점은 산술 연산의 자동 alignment 규칙을 모든 이항 연산에 그대로 적용하면 안 된다는 사실이다.

\begin{dhktable}
|연산|	Index가 다른 경우|
|:---|:---|
|a + b|	기본적으로 레이블을 맞춰 산술 연산|
|a.add(b)|	기본적으로 레이블을 맞춰 산술 연산|
|a == b|	동일하게 레이블이 배치되지 않으면 오류 가능|
|a.eq(b)|	레이블을 맞춰 비교 가능|
\end{dhktable}

DataFrame에서도 일반 비교 연산자와 \zsh{.eq()} 같은 유연한 비교 메서드 사이에 유사한 차이가 있다.

또한 \zsh{D.3}에서 배운 Nullable 자료형이라면 비교 결과에 \zsh{<NA>}가 포함될 수도 있다.

따라서 비교 연산에서는 레이블 정합성과 결측값의 비교 의미론을 모두 확인해야 한다.

\subsubsection{Series를 DataFrame에 할당할 때도 Alignment가 발생한다}

이번에는 이항 산술 연산이 아니라 할당을 살펴보자.

\begin{dhkcode}
df["prediction"] = predictions
\end{dhkcode}

이 표현식은 DataFrame에 열을 추가하거나 기존 열의 값을 변경하는 할당이다.

그런데 오른쪽에 전달하는 객체가 Series라면, Pandas는 일반적으로 Series의 Index를 DataFrame의 행 Index에 맞춘다.

\begin{dhknote}
\begin{dhksmry}[centered]
독립 실행 코드 8 - Series 할당과 NumPy 배열 할당
\end{dhksmry}
\end{dhknote}

\begin{dhkcode}
import pandas as pd

df = pd.DataFrame(
    {"label": [1, 0, 1]},
    index=["doc-A", "doc-B", "doc-C"],
)

predictions = pd.Series(
    [0, 1, 1],
    index=["doc-C", "doc-A", "doc-B"],
)

df["prediction"] = predictions

df["positional"] = predictions.to_numpy()

print(df)
\end{dhkcode}

예상 출력:

\begin{dhkzsh}
       label  prediction  positional
doc-A      1           1           0
doc-B      0           1           1
doc-C      1           0           1
\end{dhkzsh}

두 열의 값이 달라졌다.

왜 그럴까?

\begin{dhknote}
\begin{dhksmry}[centered]
Series를 할당한 경우
\end{dhksmry}
\end{dhknote}

\begin{dhkcode}
df["prediction"] = predictions
\end{dhkcode}

Pandas는 Index 레이블을 기준으로 값을 대응시킨다.

\begin{dhkzsh}
doc-A → 1
doc-B → 1
doc-C → 0
\end{dhkzsh}

\begin{dhknote}
\begin{dhksmry}[centered]
NumPy 배열을 할당한 경우
\end{dhksmry}
\end{dhknote}

\begin{dhkcode}
df["positional"] = predictions.to_numpy()
\end{dhkcode}

NumPy 배열에는 Index 레이블이 없으므로 값이 현재 행 순서에 따라 할당된다.

\begin{dhkzsh}
첫 번째 행 → 0
두 번째 행 → 1
세 번째 행 → 1
\end{dhkzsh}

여기서 중요한 사실이 있다.

Series를 할당하는 것과 Series의 값을 NumPy 배열로 변환해 할당하는 것은 서로 다른 의미의 연산이다.

\begin{dhknote}
\begin{dhksmry}[centered]
할당할 Series에 일부 레이블이 없다면?
\end{dhksmry}
\end{dhknote}

\begin{dhkcode}
import pandas as pd

df = pd.DataFrame(
    {"label": [1, 0, 1]},
    index=["doc-A", "doc-B", "doc-C"],
)

predictions = pd.Series(
    [1, 0],
    index=["doc-C", "doc-A"],
)

df["prediction"] = predictions

print(df)
\end{dhkcode}

예상 출력:

\begin{dhkzsh}
       label  prediction
doc-A      1         0.0
doc-B      0         NaN
doc-C      1         1.0
\end{dhkzsh}

\zsh{doc-B}에 대응하는 예측값이 없으므로 \zsh{NaN}이 발생한다.

또한 일반 정수형 Series에 결측이 추가되면서 결과 열의 dtype이 \zsh{float64}로 변경되었다.

이것은 \zsh{D.3}의 dtype 변환과 \zsh{D.5}의 alignment가 함께 작동한 결과다.

\begin{dhknote}
\begin{dhksmry}[centered]
객체 상태 관점에서 분석
\end{dhksmry}
\end{dhknote}

\begin{dhkcode}
df["prediction"] = predictions
\end{dhkcode}

이 표현식에서는 기존 DataFrame 객체의 열 구성이 변경된다.

반면 다음 표현식은 다르다.

\begin{dhkcode}
result = df["label"] + predictions
\end{dhkcode}

이 경우에는 새로운 Series가 반환되어 \zsh{result}에 바인딩될 뿐, \zsh{df}에 자동으로 열이 추가되지는 않는다.

즉, alignment는 값의 대응 규칙이고, 원본 변경 여부는 수행하는 연산의 종류에 따라 결정된다.

\subsubsection{중복 Index가 있으면 Alignment는 어떻게 작동할까?}

지금까지는 Index 레이블이 유일하다고 가정했다.

하지만 \zsh{D.2}에서 배웠듯이 Pandas의 Index는 중복 레이블을 허용한다.

이 경우에는 단순한 일대일 대응이 성립하지 않을 수 있다.

\begin{dhknote}
\begin{dhksmry}[centered]
독립 실행 코드 9 - 중복 레이블과 결과 크기
\end{dhksmry}
\end{dhknote}

\begin{dhkcode}
import pandas as pd

a = pd.Series(
    [1, 2],
    index=["A", "A"],
)

b = pd.Series(
    [10, 20],
    index=["A", "A"],
)

c = pd.Series(
    [10, 20, 30],
    index=["A", "A", "A"],
)

print(a + b)
print()
print(a + c)
\end{dhkcode}

예상 출력:

\begin{dhkzsh}
A    11
A    22
dtype: int64

A    11
A    21
A    31
A    12
A    22
A    32
dtype: int64
\end{dhkzsh}

첫 번째 연산에서는 두 Series의 Index가 동일하게 배치되어 있다.

따라서 같은 위치의 중복 레이블끼리 대응되어 두 개의 결과가 생성된다.

반면 두 번째 연산에서는 \zsh{"A"}가 한쪽에 두 번, 다른 쪽에 3번 등장한다.

이 경우 alignment 과정에서 \zsh{"A"}에 대응하는 값들의 조합이 만들어져 여섯 개의 결과가 발생한다.

\[
2\times3=6
\]

이것은 관계형 데이터베이스에서 중복 키를 기준으로 JOIN할 때 발생하는 다대다 (many-to-many) 결과와 유사한 현상이다.

하지만 Pandas의 산술 alignment와 SQL JOIN이 모든 세부 동작에서 동일한 것은 아니다.

여기서는 중복 키가 결과의 크기를 증가시킬 수 있다는 구조적 유사성에 주목하면 된다.

\begin{dhknote}
\begin{dhksmry}[centered]
중복 Index를 검사하는 방법
\end{dhksmry}
\end{dhknote}

\begin{dhkcode}
print(a.index.is_unique)
print(c.index.has_duplicates)
\end{dhkcode}

예상 출력:

\begin{dhkzsh}
False
True
\end{dhkzsh}

여기서 판단 기준은 명확하다.

Index가 유일하면 각 레이블은 하나의 값을 식별할 수 있다.

반면 Index가 중복되어 있으면 하나의 레이블이 여러 값을 가리킬 수 있으므로 겨로가의 개수와 대응 관계가 달라질 수 있다.

특히 모델 평가에서 문서 ID가 중복되어 있다면 예측값과 정답 레이블을 일대일로 대응시킨다는 가정 자체가 깨질 수 있다.

따라서 분석 파이프라인에서는 중요한 식별자에 대해 다음 조건을 확인하는 것이 좋다.

\begin{dhkcode}
assert df.index.is_unique
\end{dhkcode}

이것은 Index가 언제나 유일해야 한다는 뜻이 아니다.

문서 ID처럼 유일성을 요구하는 업무 규칙이 있을 때 검증해야 한다는 뜻이다.

\subsubsection{Index의 값과 자료형도 대응 관계에 영향을 준다}

레이블은 단순한 화면 표시가 아니다.

레이블 자체의 값과 타입이 대응 관계를 결정한다.

다음 예제를 살펴보자.

\begin{dhkcode}
import pandas as pd

a = pd.Series(
    [10],
    index=[1],
)

b = pd.Series(
    [5],
    index=["1"],
)

print(a + b)
\end{dhkcode}

예상 출력:

\begin{dhkzsh}
1   NaN
1   NaN
dtype: float64
\end{dhkzsh}

출력 화면에서는 두 행이 모두 1로 표시되지만, 실제 레이블은 서로 다르다.

하나는 정수 \zsh{1}이고, 다른 하나는 문자열 \zsh{"1"}이다.

따라서 Pandas는 두 레이블을 동일한 것으로 대응시키지 않는다.

다음처럼 타입을 확인할 수 있다.

\begin{dhkcode}
result = a + b

print([
    type(label).__name__ for label in result.index
])
\end{dhkcode}

예상 출력:

\begin{dhkzsh}
['int', 'str']
\end{dhkzsh}

여기서 중요한 것은 화면에 같은 문자로 출력된다고 해서 같은 레이블인 것은 아니라는 점이다.

CSV 파일을 읽거나 외부 데이터베이스에서 데이터를 가져올 때 문서 ID가 정수형과 문자열형으로 혼재하면 이러한 문제가 발생할 수 있다.

또한 \zsh{Index.name}은 축의 의미를 설명하는 메타데이터이지만, 단일 Index에서 같은 레이블의 대응 여부를 결정하는 유일한 기준은 아니다.

따라서 Index의 이름만 같다고 해서 alignment가 안전하다고 판단해서는 안 된다.

\subsubsection{Transformer 응용 - 예측 정확도가 100퍼센트에서 33.3퍼센트로 달라지는 이유}

이제 이번 단원의 개념을 하나의 실제 문제로 통합해 보자.

문서 분류 모델의 정답 레이블과 예측 레이블이 각각 Series로 저장되어 있다고 가정한다.

두 Series는 같은 문서 세 개를 포함하지만 순서가 다르다.

\begin{dhknote}
\begin{dhksmry}[centered]
독립 실행 코드 10 - 문서 ID 기반 평가
\end{dhksmry}
\end{dhknote}

\begin{dhkcode}
import pandas as pd

truth = pd.Series(
    [1, 0, 1],
    index=["doc-A", "doc-B", "doc-C"],
    dtype="Int64",
    name="label",
)

pred = pd.Series(
    [1, 0, 1],
    index=["doc-C", "doc-A", "doc-B"],
    dtype="Int64",
    name="prediction",
)

# 1. 식별자 검증
assert truth.index.is_unique
assert pred.index.is_unique

assert truth.index.difference(pred.index).empty
assert pred.index.difference(truth.index).empty

# 2. 문서 ID 기준 정렬
pred_aligned = pred.reindex(truth.index)

# 3. 레이블 기준 비교
correct = truth.eq(pred_aligned)

# 4. 정확도 계산
accuracy = correct.mean()

print(pred_aligned.tolist())
print(correct.tolist())
print(round(float(accuracy), 3))

# 5. 잘못된 위치 기반 비교
positional_correct = (
    truth.to_numpy(dtype="int64")
    == pred.to_numpy(dtype="int64")
)

print(positional_correct.tolist())
print(positional_correct.mean())
\end{dhkcode}

예상 출력:

\begin{dhkzsh}
[0, 1, 1]
[False, False, True]
0.333
[True, True, True]
1.0
\end{dhkzsh}

결과가 크게 달라졌다.

\begin{dhknote}
\begin{dhksmry}[centered]
왜 레이블 기준 정확도는 33.3%일까?
\end{dhksmry}
\end{dhknote}

정답과 예측을 문서 ID별로 정리하면 다음과 같다.

\begin{dhktable}
|문서 ID|	정답|	예측|	정답 여부|
|:---|:---|:---|:---|
|doc-A|	1	|0|	False|
|doc-B|	0	|1|	False|
|doc-C|	1	|1|	True|
\end{dhktable}

세 문서 중 하나만 맞았으므로 정확도는 다음과 같다.

\[
\text{Accuracy}=\frac{1}{3}\approx0.333
\]

반면 NumPy 배열로 변환한 뒤 위치 기준으로 비교하면 다음과 같다.

\begin{dhkzsh}
truth 값 순서: [1, 0, 1]
pred 값 순서:  [1, 0, 1]
\end{dhkzsh}

이 경우 세 위치의 값이 모두 같으므로 정확도가 100퍼센트로 계산된다.

하지만 이것은 실제 모델이 세 문서를 모두 맞혔다는 뜻이 아니다.

서로 다른 문서의 정답과 예측을 비교했기 때문에 발생한 잘못된 평가 결과다.

\begin{dhknote}
\begin{dhksmry}[centered]
왜 Index 검증을 먼저 했을까?
\end{dhksmry}
\end{dhknote}

코드 앞부분에서 다음 검사를 수행했다.

\begin{dhkcode}
assert truth.index.is_unique
assert pred.index.is_unique
\end{dhkcode}

두 Series의 문서 ID가 중복되지 않는지 확인하기 위해서다.

그리고 다음 검사를 수행했다.

\begin{dhkcode}
assert truth.index.difference(pred.index).empty
assert pred.index.difference(truth.index).empty
\end{dhkcode}

이것은 한쪽에만 존재하는 문서 ID가 없는지 확인하는 과정이다.

왜 이 검사가 필요할까?

\zsh{reindex()}는 존재하지 않는 레이블에 대해 결측값을 만들 수 있다.

따라서 누락된 예측값이 있는데도 별도의 검사 없이 정렬을 수행하면, 평가 과정에서 결측이 발생한 원인을 놓칠 수 있다.

또한 \zsh{Index.difference()}는 집합에 가까운 레이블 차이를 확인하는 것이므로 중복 개수까지 검증하는 도구는 아니다.

그래서 유일성 검사와 함께 사용하는 것이다.

이 구조는 모델 평가 파이프라인에서 다음 순서를 갖는다는 점이 중요하다.

\begin{dhksmry}[centered]
정답 Series / 예측 Series
각 객체가 문서 ID를 Index로 보유
\da
Index 유일성 및 레이블 집합 검증
중복·누락·불필요한 예측 ID 검사
\da
reindex() 또는 align()
문서 ID 기준으로 값의 대응 관계 확정
\da
정답 비교 및 평가 지표 계산
Accuracy · F1 · 오류 샘플 분석
\end{dhksmry}

이 과정에서 Index는 단순한 행 이름이 아니라 서로 다른 데이터 객체의 정합성을 유지하는 식별 체계로 작동한다.

이것은 이후 Transformer의 토크노하 결과, \zsh{Dataset}, \zsh{DataLoader}, 예측 결과를 Pandas로 다시 결합할 때도 적용되는 원칙이다.

단, PyTorch Tensor 자체는 Pandas Index를 보유하지 않는다.

따라서 Tensor로 변환하는 순간부터는 샘플 ID와 행 순서의 대응을 별도로 관리해야 한다.

이 문제는 \zsh{D.21}에서 더 깊게 다룰 것이다.

\subsubsection{기존 Python 개념과 연결하기}

이번 학습의 핵심은 Pandas가 Python의 일반적인 객체 모델을 기반으로 레이블 정합이라는 추가적인 연산 규칙을 구현했다는 사실이다.

\begin{dhktable}
|Python 객체 모델 개념|	Pandas에서의 적용|
|:---|:---|
|객체 생성|	Series·DataFrame·Index 인스턴스 생성|
|이름 바인딩|	a, b, result가 각각 객체를 참조|
|연산자 오버로딩|	a + b가 Pandas의 산술 연산 규칙을 사용|
|속성 탐색|	a.index, a.dtype 등 조회|
|메서드 호출|	a.add(b), a.align(b), a.reindex(...)|
|반환 객체|	산술 연산은 새로운 Series/DataFrame을 반환|
|상태 변경|	df["prediction"] = s는 기존 DataFrame에 값을 할당|
|타입과 자료형|	type(s)와 s.dtype은 다른 계층의 정보|
|동등성|	Index.equals()는 레이블 구조 비교에 사용|
|객체 동일성|	is는 동일한 Python 객체인지 비교할 뿐 레이블 정합을 판단하지 않음|
\end{dhktable}

여기서 특히 구분해야 할 개념이 있다.

\begin{dhkcode}
a.index is b.index
\end{dhkcode}

이 표현식은 두 Series가 정확히 같은 Index 객체를 참조하는지 검사한다.

반면 다음 표현식은 다르다.

\begin{dhkcode}
a.index.equals(b.index)
\end{dhkcode}

이것은 두 Index의 레이블과 순서가 같은지 비교하는 메서드다.

두 Series가 서로 다른 Index 객체를 갖고 있더라도 레이블이 같다면 alignment 관점에서는 대응될 수 있다.

따라서 Index 객체의 동일성과 Index 레이블의 정합성은 서로 다른 개념이다.

또한 \zsh{Index.equals()} 만으로 데이터 품질의 모든 조건이 검증되는 것은 아니다.

레이블의 유일성, 누락 여부, 업무상 식별자의 의미까지 필요에 따라 별도로 검사해야 한다.

\subsubsection{핵심 정리와 판단 기준}

이번 학습에서 가장 중요한 사실은 네 가지다.

첫째, Series 간 기본 산술 연산은 위치가 아니라 레이블을 기준으로 값을 대응시킨다.

따라서 같은 \zsh{shape}을 가진 두 Series라도 Index의 순서나 레이블 집합이 다르면 결과가 달라질 수 있다.

둘째, DataFrame 간 산술 연산에서는 행과 열의 레이블이 모두 개입한다.

두 DataFrame의 \zsh{shape}이 같더라도 연산 결과의 \zsh{shape}은 달라질 수 있다.

셋째, 연산 종류에 따라 Alignment 규칙이 달라진다.

산술 연산, 비교 연산, Series 할당, NumPy 배열 할당은 모두 동일한 규칙을 사용하는 것이 아니다.

넷째, Alignment는 데이터 정합성을 유지하는 데 유용하지만, 데이터의 올바름을 자동으로 보장하지는 않는다.

중복 ID, 누락 ID, 잘못된 레이블 자료형이 있다면 alignment가 수행되더라도 결과가 잘못되거나 예상보다 커질 수 있다.

\begin{dhknote}
\begin{dhksmry}[centered]
실전 판단 기준
\end{dhksmry}
\end{dhknote}

\begin{dhktable}
|확인할 질문|	판단 기준|
|:---|:---|
|두 객체의 크기가 같은가?|	shape 확인|
|같은 원소를 식별하는가?	|Index 레이블의 의미 확인|
|레이블 집합이 같은가?	|Index.difference() 등으로 확인|
|레이블 순서까지 같은가?|	Index.equals() 확인|
|레이블이 유일한가?	|Index.is_unique 확인|
|어떤 축을 기준으로 연산하는가?|	axis 확인|
|결측이 발생할 수 있는가?	|Index 합집합과 dtype 확인|
|원본이 변경되는가?|	산술 연산인지 할당인지 구분|
|NumPy로 변환해도 안전한가?|	레이블 제거 후 행 순서가 보장되는지 확인|
\end{dhktable}

이 기준을 익히면 Pandas의 연산 결과를 단순히 실행해 보는 수준을 넘어, 실행 전에 객체와 Index의 관계를 근거로 예측할 수 있게 된다.

\subsubsection{다음 학습}

이번 \zsh{D.5}에서는 Index가 산술 연산과 할당 과정에 어떻게 개입하는지 배웠다.

다음 학습에서는 이러한 Index alignment를 불리언 마스크와 연결할 것이다.

\begin{dhkcode}
mask = df["score"] >= 80

result = df.loc[mask]
\end{dhkcode}

여기서 \zsh{mask}는 어떤 객체일까?

\zsh{df.loc[mask]}는 마스크의 위치를 기준으로 행을 선택할까, 아니면 Index 레이블을 기준으로 선택할까?

그리고 마스크에 \zsh{<NA>}가 포함되어 있다면 어떻게 될까?

\zsh{D.6}에서는 \zsh{isin()}, \zsh{where()}, \zsh{mask()}, \zsh{query()}까지 확장하면서 조건을 계산하는 과정과 조건을 적용하는 과정을 구분해 분석할 것이다.

공식 참고 자료

\begin{dhksmry}
https://pandas.pydata.org/docs/user_guide/basics.html
https://pandas.pydata.org/docs/user_guide/indexing.html
https://pandas.pydata.org/docs/reference/api/pandas.Series.align.html
\end{dhksmry}